% Save as: module2-problem-set.tex
% ============================================================================
% Problem Set 2: Consumer Theory, Demand & Welfare
% Microeconomic Theory for Experimentalists & Applied Microeconomists
% Built from templates/problem_set_template.tex. Compile with pdflatex.
% ============================================================================
\documentclass[11pt]{article}
\usepackage{amsmath, amssymb, amsthm}
\usepackage[margin=1in]{geometry}
\usepackage{enumitem}

\newtheorem{question}{Question}

\title{Problem Set 2: Consumer Theory, Demand \& Welfare\\
\large Microeconomic Theory for Experimentalists \& Applied Microeconomists}
\author{Due: \textbf{[date --- before Lecture 1 of Module 3]}}
\date{}

\begin{document}
\maketitle

\noindent\emph{Ground rules: collaboration on ideas is encouraged; write-ups
are individual. You may use AI tools for parts flagged with
$\langle$AI$\rangle$ and must attach the transcript. Show all calculations.}

\medskip

\noindent\textbf{Weights:} Part A --- 30 points. Part B --- 45 points.
Part C --- 25 points.

\bigskip

% ---------------------------------------------------------------------------
\section*{Part A: Theory --- Guided Calculations (30 points)}

\begin{question}[Demand and the Slutsky decomposition --- 15 points]
A consumer has utility $u(x, y) = xy$, income $m = 24$, and faces prices
$p_y = 3$ throughout.

\begin{enumerate}[label=(\alph*)]
  \item \textbf{(4 pts)} At the optimum, the marginal rate of substitution
        $y/x$ equals the price ratio $p_x/p_y$. Use this and the budget
        constraint to show that demand is $x = \dfrac{m}{2 p_x}$ and
        $y = \dfrac{m}{2 p_y}$ --- i.e.\ she spends half her income on
        each good. Compute the bundle at $p_x = 3$.
  \item \textbf{(3 pts)} The price of $x$ falls to $p_x = 2$. Compute the
        new bundle and state the total effect on $x$.
  \item \textbf{(5 pts)} Slutsky decomposition. Compute the compensated
        income $m'$ that makes the \emph{original} bundle exactly
        affordable at the new prices. Compute demand for $x$ at
        $(p_x = 2,\; m')$, and report the substitution and income effects
        separately. \emph{Hint: $m' = m - x_{\text{old}} \cdot \Delta p_x$
        where $\Delta p_x$ is the price DROP; both effects here are whole
        numbers.}
  \item \textbf{(3 pts)} A development agency is choosing between a price
        subsidy on good $x$ and an equivalent-cost cash transfer (equal
        total budgetary outlay). In one
        paragraph, use your decomposition to explain what an experiment
        that measures only the \emph{total} demand response to the subsidy
        cannot tell the agency, and what treatment arm fixes it.
\end{enumerate}
\end{question}

\begin{question}[CV, EV, and consumer surplus with numbers --- 15 points]
A consumer has quasi-linear utility: her marginal value for good $x$ is
$20 - 2x$ dollars, and utility is linear in money spent on everything else
(so there are no income effects on $x$). The price of $x$ rises from
\$4 to \$8.

\begin{enumerate}[label=(\alph*)]
  \item \textbf{(3 pts)} Compute her demand for $x$ at each price.
        \emph{Hint: she buys until marginal value equals price.}
  \item \textbf{(5 pts)} Compute the loss in consumer surplus from the
        price rise (a trapezoid --- draw it). Explain in one sentence why,
        with quasi-linear utility, CV and EV both equal this number.
  \item \textbf{(4 pts)} Now suppose $x$ were instead a \emph{normal good}
        with income effects. State which is larger \emph{in absolute
        value} for this price rise --- CV or EV --- and give the intuition
        in two sentences. \emph{Hint:
        CV asks the compensation question at NEW prices, when the consumer
        is poorer in real terms.}
  \item \textbf{(3 pts)} A ministry must budget compensation payments for
        households hurt by this price rise, and separately wants to
        predict how households would vote on a referendum to block it.
        Which welfare measure belongs to which task, and why?
\end{enumerate}
\end{question}

% ---------------------------------------------------------------------------
\section*{Part B: Experimental Design (45 points)}

\begin{question}[Is the WTP--WTA gap an endowment effect?]
Kahneman, Knetsch \& Thaler (1990) found median WTA roughly twice median
WTP for the same mug. The loss-aversion reading says this is a real
feature of preferences. The skeptical reading (Plott--Zeisel style) says
much of it is subject confusion about the elicitation mechanism and
strategic habits imported from everyday bargaining. Design a lab
experiment that measures the WTP--WTA gap AND separates these two
readings. Specify:
\begin{enumerate}[label=(\alph*)]
  \item \textbf{(6 pts)} Subject pool and recruitment, and why prior
        experience with the mechanism (or its absence) matters for your
        question.
  \item \textbf{(12 pts)} Treatments. Your design must include (i) an
        endowment manipulation (sellers vs.\ buyers of the same good) and
        (ii) at least one arm that changes subjects' \emph{understanding
        or experience} of the mechanism (e.g.\ paid practice rounds on
        induced-value tokens) without changing the good or the endowment.
        Say what each comparison identifies.
  \item \textbf{(9 pts)} The elicitation mechanism and its
        incentive-compatibility argument, written for a skeptic in plain
        language (this is the tutor-session drill). Then name one failure
        mode the mechanism does NOT fix.
  \item \textbf{(9 pts)} Hypotheses and rough sample-size reasoning:
        state the WTA/WTP ratio you expect in each arm under each reading
        (anchor on KKT's $\approx 2$ and on the training-shrinks-the-gap
        literature), and the approximate $n$ per arm to distinguish them.
        Simulation-based reasoning is welcome; formulas are not required.
  \item \textbf{(9 pts)} The biggest confound remaining in YOUR design and
        your answer to it.
\end{enumerate}
\end{question}

% ---------------------------------------------------------------------------
\section*{Part C: Silicon Pilot Interpretation $\langle$AI$\rangle$
(25 points)}

\begin{question}[The agents that were too rational]
You run \texttt{module2-simulation-lab.ipynb}. Across all personas and
framings, your agents' four-menu choices come back \textbf{100\%
WARP-consistent} --- not a single violation in 180 responses. Human
subjects in comparable budget-choice experiments (Andreoni \& Miller 2002;
Choi et al.\ 2007) are highly but \emph{imperfectly} consistent: roughly
one in ten violates GARP, and consistency scores vary meaningfully across
people. (The comparison is qualitative: the lab's audit is pairwise WARP
over four menus, while the human statistics are GARP rates over longer
menu sequences.) Your lab partner reports: ``the agents replicate rational
consumer theory --- great news.''

\begin{enumerate}[label=(\alph*)]
  \item \textbf{(8 pts)} Explain why perfect consistency is a
        \emph{divergence} from human behavior, not a replication ---
        and give two mechanisms (at least one LLM-specific) that could
        produce it.
  \item \textbf{(9 pts)} Design the diagnostic simulation that tests
        whether silicon consistency is robust or fragile: propose at least
        two manipulations (e.g.\ non-round prices, more menus, distractor
        text in the prompt, a persona described as busy or distracted) and
        state what result pattern under each manipulation would tell you.
        Run it if API budget allows and report; otherwise state the
        predicted patterns precisely.
  \item \textbf{(8 pts)} Your research group wants to use this lab to
        pilot a budget-choice experiment for human subjects --- including
        choosing the number of menus and estimating how many subjects
        will fail rationality screens. State one thing the silicon pilot
        CAN safely inform, one thing it CANNOT (and why), and the design
        change this implies for the human pilot plan.
\end{enumerate}
Attach your AI transcripts and, if you ran code, the results CSV.
\end{question}

\end{document}
